\GalSourcePageBoundary{072}{L0071}{43}{43}

Adjoignons à l'équation le premier radical extrait dans la solution.
Il pourra arriver deux cas: ou bien, par l'adjonction de ce
radical, le groupe des permutations de l'équation sera diminué; ou
bien, cette extraction de racine n'étant qu'une simple préparation,
le groupe restera le même.

Toujours sera-t-il qu'après un certain nombre \emph{fini} d'extractions
de racines, le groupe devra se trouver diminué, sans quoi l'équation
ne serait pas soluble.

Si, arrivé à ce point, il y avait plusieurs manières de diminuer le
groupe de l'équation proposée par une simple extraction de racine,
il faudrait, pour ce que nous allons dire, considérer seulement un
radical du degré le moins haut possible parmi tous les simples
radicaux, qui sont tels que la connaissance de chacun d'eux
diminue le groupe de l'équation.

Soit donc $p$ le nombre premier qui représente ce degré minimum,
en sorte que par une extraction de racine de degré $p$, on
diminue le groupe de l'équation.

Nous pouvons toujours supposer, du moins pour ce qui est
relatif au groupe de l'équation, que, parmi les quantités adjointes
précédemment à l'équation, se trouve une racine $p^{\text{ième}}$ de l'unité, $\alpha$.
Car, comme cette expression s'obtient par des extractions de
racines de degré inférieur à $p$, sa connaissance n'altérera en rien
le groupe de l'équation.

Par conséquent, d'après les théorèmes II et III, le groupe de
l'équation devra se décomposer en $p$ groupes jouissant les uns par
rapport aux autres de cette double propriété: 1\textsuperscript{o} que l'on passe
de l'un à l'autre par une seule et même substitution; 2\textsuperscript{o} que tous
contiennent les mêmes substitutions.

Je dis réciproquement que, si le groupe de l'équation peut se
partager en $p$ groupes qui jouissent de cette double propriété, on
pourra, par une simple extraction de racine $p^{\text{ième}}$, et par l'adjonction
de cette racine $p^{\text{ième}}$, réduire le groupe de l'équation à l'un de
ces groupes partiels.

Prenons, en effet, une fonction des racines qui soit invariable
pour toutes les substitutions de l'un des groupes partiels, et varie
pour toute autre substitution. (Il suffit, pour cela, de choisir une
fonction symétrique des diverses valeurs que prend, par toutes les
